\documentclass{amsart}
% For dealing with references we use the comment environment
\usepackage{verbatim}
\newenvironment{reference}{\comment}{\endcomment}
%\newenvironment{reference}{}{}
\newenvironment{slogan}{\comment}{\endcomment}
\newenvironment{history}{\comment}{\endcomment}

% For commutative diagrams we use Xy-pic
\usepackage[all]{xy}

% We use 2cell for 2-commutative diagrams.
\xyoption{2cell}
\UseAllTwocells

% We use multicol for the list of chapters between chapters
\usepackage{multicol}

% This is generally recommended for better output
\usepackage{lmodern}
\usepackage[T1]{fontenc}

% For cross-file-references

% Package for hypertext links:
\usepackage{hyperref}

% For any local file, say "hello.tex" you want to link to please

% Theorem environments.
%
\theoremstyle{plain}
\newtheorem{theorem}[subsection]{Theorem}
\newtheorem{proposition}[subsection]{Proposition}
\newtheorem{lemma}[subsection]{Lemma}

\theoremstyle{definition}
\newtheorem{definition}[subsection]{Definition}
\newtheorem{example}[subsection]{Example}
\newtheorem{exercise}[subsection]{Exercise}
\newtheorem{situation}[subsection]{Situation}

\theoremstyle{remark}
\newtheorem{remark}[subsection]{Remark}
\newtheorem{remarks}[subsection]{Remarks}

\numberwithin{equation}{subsection}

% Macros
%
\def\lim{\mathop{\mathrm{lim}}\nolimits}
\def\colim{\mathop{\mathrm{colim}}\nolimits}
\def\Spec{\mathop{\mathrm{Spec}}}
\def\Hom{\mathop{\mathrm{Hom}}\nolimits}
\def\Ext{\mathop{\mathrm{Ext}}\nolimits}
\def\SheafHom{\mathop{\mathcal{H}\!\mathit{om}}\nolimits}
\def\SheafExt{\mathop{\mathcal{E}\!\mathit{xt}}\nolimits}
\def\Sch{\mathit{Sch}}
\def\Mor{\mathop{\mathrm{Mor}}\nolimits}
\def\Ob{\mathop{\mathrm{Ob}}\nolimits}
\def\Sh{\mathop{\mathit{Sh}}\nolimits}
\def\NL{\mathop{N\!L}\nolimits}
\def\CH{\mathop{\mathrm{CH}}\nolimits}
\def\proetale{{pro\text{-}\acute{e}tale}}
\def\etale{{\acute{e}tale}}
\def\QCoh{\mathit{QCoh}}
\def\Ker{\mathop{\mathrm{Ker}}}
\def\Im{\mathop{\mathrm{Im}}}
\def\Coker{\mathop{\mathrm{Coker}}}
\def\Coim{\mathop{\mathrm{Coim}}}

% Boxtimes
%
\DeclareMathSymbol{\boxtimes}{\mathbin}{AMSa}{"02}

%
% Macros for moduli stacks/spaces
%
\def\QCohstack{\mathcal{QC}\!\mathit{oh}}
\def\Cohstack{\mathcal{C}\!\mathit{oh}}
\def\Spacesstack{\mathcal{S}\!\mathit{paces}}
\def\Quotfunctor{\mathrm{Quot}}
\def\Hilbfunctor{\mathrm{Hilb}}
\def\Curvesstack{\mathcal{C}\!\mathit{urves}}
\def\Polarizedstack{\mathcal{P}\!\mathit{olarized}}
\def\Complexesstack{\mathcal{C}\!\mathit{omplexes}}
% \Pic is the operator that assigns to X its picard group, usage \Pic(X)
% \Picardstack_{X/B} denotes the Picard stack of X over B
% \Picardfunctor_{X/B} denotes the Picard functor of X over B
\def\Pic{\mathop{\mathrm{Pic}}\nolimits}
\def\Picardstack{\mathcal{P}\!\mathit{ic}}
\def\Picardfunctor{\mathrm{Pic}}
\def\Deformationcategory{\mathcal{D}\!\mathit{ef}}

\setlength{\emergencystretch}{3em}
\begin{document}
\title[Stacks Project, Tag 0C1F]{569 --- Riemann\textendash{}Hurwitz multiplicities and the tame equality criterion for generically étale smooth proper curves with ground-field global functions}
\maketitle
\noindent Copyright (C) 2005--2025 Johan de Jong. Stacks Project authors. Source: \href{https://stacks.math.columbia.edu/tag/0C1F}{Tag 0C1F}.

\noindent Editorial difficulty note (not author text). Difficulty H2: The proof must treat inseparable residue fields uniformly rather than only differentiate a separable local uniformizer. Base change of the Gorenstein different reduces the calculation to an Artinian fibre algebra, where the trace element detects tame ramification and its annihilator determines the exact different ideal..

\noindent Editorial provenance note (not author text). History: excerpt from revision \texttt{a0444\allowbreak 6e57e\allowbreak c1fbc\allowbreak 252a8\allowbreak 71afc\allowbreak ec775\allowbreak 2fb28\allowbreak 07b14}, curves.tex, lines 2432--2523. Reference presentation was adapted; original-author context and core support blocks were retained or added. The mathematical wording is unchanged. Licensed under GNU FDL 1.2 or later; no invariant sections or cover texts. See ../licenses/COPYING. Context blocks reproduce author definitions and supporting results; they are not additional counted problems.

\begin{definition}
\label{discriminant-definition-different}
Let $f : Y \to X$ be a flat locally quasi-finite morphism of
locally Noetherian schemes.
Let $\omega_{Y/X}$ be the relative dualizing module and let
$\tau_{Y/X} \in \Gamma(Y, \omega_{Y/X})$ be the trace element
(Remarks \href{https://stacks.math.columbia.edu/tag/0BVG}{remark relative dualizing for quasi finite (Tag 0BVG)} and
\href{https://stacks.math.columbia.edu/tag/0BVJ}{remark relative dualizing for flat quasi finite (Tag 0BVJ)}).
The annihilator of
$$
\Coker(\mathcal{O}_Y \xrightarrow{\tau_{Y/X}} \omega_{Y/X})
$$
is the {\it different} of $Y/X$. It is a coherent ideal
$\mathfrak{D}_f \subset \mathcal{O}_Y$.
\end{definition}

\begin{lemma}
\label{lemma-rh}
Let $f : X \to Y$ be a morphism of smooth proper curves
over a field $k$ which satisfies the equivalent conditions of
Lemma \href{https://stacks.math.columbia.edu/tag/0C1C}{lemma generically etale (Tag 0C1C)}. If
$k = H^0(Y, \mathcal{O}_Y) = H^0(X, \mathcal{O}_X)$
and $X$ and $Y$ have genus $g_X$ and $g_Y$, then
$$
2g_X - 2 = (2g_Y - 2) \deg(f) + \deg(R)
$$
where $R \subset X$ is the effective Cartier divisor cut out by
the different of $f$.
\end{lemma}

\section{Riemann-Hurwitz}
\label{section-riemann-hurewitz}

\noindent
Let $k$ be a field. Let $f : X \to Y$ be a morphism of smooth curves over $k$.
Then we obtain a canonical exact sequence
$$
f^*\Omega_{Y/k} \xrightarrow{\text{d}f} \Omega_{X/k}
\longrightarrow \Omega_{X/Y} \longrightarrow 0
$$
by Morphisms, Lemma \href{https://stacks.math.columbia.edu/tag/01UX}{lemma triangle differentials (Tag 01UX)}.
Since $X$ and $Y$ are smooth, the sheaves $\Omega_{X/k}$ and
$\Omega_{Y/k}$ are invertible modules, see
Morphisms, Lemma \href{https://stacks.math.columbia.edu/tag/02G1}{lemma smooth omega finite locally free (Tag 02G1)}.
Assume the first map is nonzero, i.e., assume $f$ is generically
\'etale, see Lemma \href{https://stacks.math.columbia.edu/tag/0C1C}{lemma generically etale (Tag 0C1C)}. Let $R \subset X$
be the closed subscheme cut out by the different $\mathfrak{D}_f$ of $f$.
By Discriminants, Lemma
\href{https://stacks.math.columbia.edu/tag/0BWJ}{lemma discriminant quasi finite morphism smooth (Tag 0BWJ)}
this is the same as the vanishing locus of $\text{d}f$, it is
an effective Cartier divisor, and we get
$$
f^*\Omega_{Y/k} \otimes_{\mathcal{O}_X} \mathcal{O}_X(R) = \Omega_{X/k}
$$
In particular, if $X$, $Y$ are projective with
$k = H^0(Y, \mathcal{O}_Y) = H^0(X, \mathcal{O}_X)$
and $X$, $Y$ have genus $g_X$, $g_Y$, then we get the
Riemann-Hurwitz formula
\begin{align*}
2g_X - 2 & =
\deg(\Omega_{X/k}) \\
& =
\deg(f^*\Omega_{Y/k} \otimes_{\mathcal{O}_X} \mathcal{O}_X(R)) \\
& =
\deg(f) \deg(\Omega_{Y/k}) + \deg(R) \\
& =
\deg(f) (2g_Y - 2) + \deg(R)
\end{align*}
The first and last equality by Lemma \href{https://stacks.math.columbia.edu/tag/0C1A}{lemma genus smooth (Tag 0C1A)}.
The second equality by the isomorphism of invertible sheaves given above.
The third equality by additivity of degrees
(Varieties, Lemma \href{https://stacks.math.columbia.edu/tag/0AYX}{lemma degree tensor product (Tag 0AYX)}),
the formula for the degree of a pullback
(Varieties, Lemma \href{https://stacks.math.columbia.edu/tag/0AYZ}{lemma degree pullback map proper curves (Tag 0AYZ)}),
and finally the formula for the degree of $\mathcal{O}_X(R)$
(Varieties, Lemma \href{https://stacks.math.columbia.edu/tag/0AYY}{lemma degree effective Cartier divisor (Tag 0AYY)}).

\medskip\noindent
To use the Riemann-Hurwitz formula we need to compute
$\deg(R) = \dim_k \Gamma(R, \mathcal{O}_R)$. By the structure
of zero dimensional schemes over $k$ (see for example
Varieties, Lemma \href{https://stacks.math.columbia.edu/tag/06LH}{lemma algebraic scheme dim 0 (Tag 06LH)}),
we see that $R$ is a finite disjoint union of spectra of
Artinian local rings $R = \coprod_{x \in R} \Spec(\mathcal{O}_{R, x})$
with each $\mathcal{O}_{R, x}$ of finite dimension over $k$. Thus
$$
\deg(R) = \sum\nolimits_{x \in R} \dim_k \mathcal{O}_{R, x} =
\sum\nolimits_{x \in R} d_x [\kappa(x) : k]
$$
with
$$
d_x = \text{length}_{\mathcal{O}_{R, x}} \mathcal{O}_{R, x} =
\text{length}_{\mathcal{O}_{X, x}} \mathcal{O}_{R, x}
$$
the multiplicity of $x$ in $R$
(see Algebra, Lemma \href{https://stacks.math.columbia.edu/tag/02M0}{lemma pushdown module (Tag 02M0)}).
Let $x \in X$ be a closed point with image $y \in Y$.
Looking at stalks we obtain an exact sequence
$$
\Omega_{Y/k, y} \to \Omega_{X/k, x} \to \Omega_{X/Y, x} \to 0
$$
Choosing local generators $\eta_x$ and $\eta_y$ of the
(free rank $1$) modules $\Omega_{X/k, x}$ and $\Omega_{Y/k, y}$
we see that
$
\eta_y \mapsto h \eta_x
$
for some nonzero $h \in \mathcal{O}_{X, x}$. By the exact sequence we see that
$\Omega_{X/Y, x} \cong \mathcal{O}_{X, x}/h\mathcal{O}_{X, x}$
as $\mathcal{O}_{X, x}$-modules. Since the divisor
$R$ is cut out by $h$ (see above) we have
$\mathcal{O}_{R, x} = \mathcal{O}_{X, x}/h\mathcal{O}_{X, x}$.
Thus we find the following equalities
\begin{align*}
d_x
& =
\text{length}_{\mathcal{O}_{X, x}}(\mathcal{O}_{R, x}) \\
& =
\text{length}_{\mathcal{O}_{X, x}}(\mathcal{O}_{X, x}/h\mathcal{O}_{X, x}) \\
& =
\text{length}_{\mathcal{O}_{X, x}}(\Omega_{X/Y, x}) \\
& =
\text{ord}_{\mathcal{O}_{X, x}}(h) \\
& =
\text{ord}_{\mathcal{O}_{X, x}}(``\eta_y/\eta_x")
\end{align*}
The first equality by our definition of $d_x$. The second and third
we saw above. The fourth equality is the definition of $\text{ord}$, see
Algebra, Definition \ref{algebra-definition-ord}. Note that since
$\mathcal{O}_{X, x}$ is a discrete valuation ring, the integer
$\text{ord}_{\mathcal{O}_{X, x}}(h)$ just the valuation of $h$.
The fifth equality is a mnemonic.

\medskip\noindent
Here is a case where one can ``calculate'' the multiplicity $d_x$ in terms
of other invariants. Namely, if $\kappa(x)$ is separable over $k$, then
we may choose $\eta_x = \text{d}s$ and $\eta_y = \text{d}t$ where $s$
and $t$ are uniformizers in $\mathcal{O}_{X, x}$ and $\mathcal{O}_{Y, y}$
(Lemma \href{https://stacks.math.columbia.edu/tag/0C1E}{lemma uniformizer works (Tag 0C1E)}).
Then $t \mapsto u s^{e_x}$ for some unit $u \in \mathcal{O}_{X, x}$
where $e_x$ is the ramification index of the extension
$\mathcal{O}_{Y, y} \subset \mathcal{O}_{X, x}$. Hence we get
$$
\eta_y = \text{d}t = \text{d}(u s^{e_x}) =
e s^{e_x - 1} u \text{d}s + s^{e_x} \text{d}u
$$
Writing $\text{d}u = w \text{d}s$ for some $w \in \mathcal{O}_{X, x}$
we see that
$$
``\eta_y/\eta_x" = e s^{e_x - 1} u + s^{e_x} w = (e_x u + s w)s^{e_x - 1}
$$
We conclude that the order of vanishing of this is $e_x - 1$
unless the characteristic of $\kappa(x)$ is $p > 0$ and $p$ divides $e_x$
in which case the order of vanishing is $> e_x - 1$.

\medskip\noindent
Combining all of the above we find that if $k$ has characteristic
zero, then
$$
2g_X - 2 = (2g_Y - 2)\deg(f) +
\sum\nolimits_{x \in X} (e_x - 1)[\kappa(x) : k]
$$
where $e_x$ is the ramification index of $\mathcal{O}_{X, x}$ over
$\mathcal{O}_{Y, f(x)}$. This precise formula will hold if and only
if all the ramification is tame, i.e., when the
residue field extensions $\kappa(x)/\kappa(y)$ are separable and
$e_x$ is prime to the characteristic of $k$, although the
arguments above are insufficient to prove this. We refer the reader
to Lemma \ref{lemma-rhe} and its proof.



\begin{definition}
\label{algebra-definition-ord}
Suppose that $K$ is a field, and $R \subset K$ is a
local\footnote{We could also define this when $R$ is only
semi-local but this is probably never really what you want!}
Noetherian subring of dimension $1$ with fraction field $K$.
In this case we define the {\it order of vanishing along $R$}
$$
\text{ord}_R : K^* \longrightarrow \mathbf{Z}
$$
by the rule
$$
\text{ord}_R(x) = \text{length}_R(R/(x))
$$
if $x \in R$ and we set
$\text{ord}_R(x/y) = \text{ord}_R(x) - \text{ord}_R(y)$
for $x, y \in R$ both nonzero.
\end{definition}

\begin{lemma}
\label{lemma-rhe}
Notation and assumptions as in Lemma \ref{lemma-rh}. For a closed point
$x \in X$ let $d_x$ be the multiplicity of $x$ in $R$. Then
$$
2g_X - 2 = (2g_Y - 2) \deg(f) + \sum\nolimits d_x [\kappa(x) : k]
$$
Moreover, we have the following results
\begin{enumerate}
\item $d_x = \text{length}_{\mathcal{O}_{X, x}}(\Omega_{X/Y, x})$,
\item $d_x \geq e_x - 1$ where $e_x$ is the ramification index
of $\mathcal{O}_{X, x}$ over $\mathcal{O}_{Y, y}$,
\item $d_x = e_x - 1$ if and only if $\mathcal{O}_{X, x}$ is tamely
ramified over $\mathcal{O}_{Y, y}$.
\end{enumerate}
\end{lemma}

\begin{proof}
By Lemma \ref{lemma-rh} and the discussion above
(which used
Varieties, Lemma \href{https://stacks.math.columbia.edu/tag/06LH}{lemma algebraic scheme dim 0 (Tag 06LH)}
and
Algebra, Lemma \href{https://stacks.math.columbia.edu/tag/02M0}{lemma pushdown module (Tag 02M0)})
it suffices to prove the results on the
multiplicity $d_x$ of $x$ in $R$. Part (1) was proved
in the discussion above. In the discussion above
we proved (2) and (3) only in the case where $\kappa(x)$ is
separable over $k$.
In the rest of the proof we give a uniform treatment
of (2) and (3) using material on differents of
quasi-finite Gorenstein morphisms.

\medskip\noindent
First, observe that $f$ is a quasi-finite Gorenstein morphism.
This is true for example because
$f$ is a flat quasi-finite morphism and $X$ is Gorenstein
(see Duality for Schemes, Lemma
\href{https://stacks.math.columbia.edu/tag/0C12}{lemma flat morphism from gorenstein scheme (Tag 0C12)})
or because it was shown in the proof of
Discriminants, Lemma
\href{https://stacks.math.columbia.edu/tag/0BWJ}{lemma discriminant quasi finite morphism smooth (Tag 0BWJ)}
(which we used above). Thus $\omega_{X/Y}$ is invertible by
Discriminants, Lemma \href{https://stacks.math.columbia.edu/tag/0C16}{lemma gorenstein quasi finite (Tag 0C16)}
and the same remains true after replacing $X$ by opens and after
performing a base change by some $Y' \to Y$. We will use this
below without further mention.

\medskip\noindent
Choose affine opens $U \subset X$ and $V \subset Y$
such that $x \in U$, $y \in V$, $f(U) \subset V$, and $x$ is the only
point of $U$ lying over $y$. Write $U = \Spec(A)$ and $V = \Spec(B)$.
Then $R \cap U$ is the different of $f|_U : U \to V$.
By Discriminants, Lemma \href{https://stacks.math.columbia.edu/tag/0BW7}{lemma base change different (Tag 0BW7)}
formation of the different commutes with arbitrary base change
in our case. By our choice of $U$ and $V$ we have
$$
A \otimes_B \kappa(y) =
\mathcal{O}_{X, x} \otimes_{\mathcal{O}_{Y, y}} \kappa(y) =
\mathcal{O}_{X, x}/(s^{e_x})
$$
where $e_x$ is the ramification index as in the statement of the lemma.
Let $C = \mathcal{O}_{X, x}/(s^{e_x})$ viewed as a finite algebra
over $\kappa(y)$. Let $\mathfrak{D}_{C/\kappa(y)}$ be the different
of $C$ over $\kappa(y)$ in the sense of
Discriminants, Definition \ref{discriminant-definition-different}.
It suffices to show: $\mathfrak{D}_{C/\kappa(y)}$
is nonzero if and only if the extension
$\mathcal{O}_{Y, y} \subset \mathcal{O}_{X, x}$ is tamely ramified
and in the tamely ramified case $\mathfrak{D}_{C/\kappa(y)}$
is equal to the ideal generated by $s^{e_x - 1}$ in $C$.
Recall that tame ramification means exactly that $\kappa(x)/\kappa(y)$
is separable and that the characteristic of $\kappa(y)$ does not
divide $e_x$. On the other hand, the different of $C/\kappa(y)$ is nonzero
if and only if $\tau_{C/\kappa(y)} \in \omega_{C/\kappa(y)}$ is nonzero.
Namely, since $\omega_{C/\kappa(y)}$ is an invertible $C$-module
(as the base change of $\omega_{A/B}$)
it is free of rank $1$, say with generator $\lambda$. Write
$\tau_{C/\kappa(y)} = h\lambda$ for some $h \in C$. Then
$\mathfrak{D}_{C/\kappa(y)} = (h) \subset C$ whence the claim.
By Discriminants, Lemma \href{https://stacks.math.columbia.edu/tag/0C13}{lemma tau nonzero (Tag 0C13)}
we have $\tau_{C/\kappa(y)} \not = 0$
if and only if $\kappa(x)/\kappa(y)$
is separable and $e_x$ is prime to the characteristic.
Finally, even if $\tau_{C/\kappa(y)}$ is nonzero, then
it is still the case that $s \tau_{C/\kappa(y)} = 0$
because $s\tau_{C/\kappa(y)} : C \to \kappa(y)$
sends $c$ to the trace of the nilpotent operator $sc$ which is zero.
Hence $sh = 0$, hence $h \in (s^{e_x - 1})$ which proves
that $\mathfrak{D}_{C/\kappa(y)} \subset (s^{e_x - 1})$ always.
Since $(s^{e_x - 1}) \subset C$ is the smallest nonzero ideal,
we have proved the final assertion.
\end{proof}
\end{document}
